\documentclass[12pt]{article}

\usepackage{amsmath, amssymb}
\usepackage{geometry}
\usepackage{xcolor}
\pagecolor{black}
\color{white}
\geometry{margin=1in}

\title{Vector Analysis}
\author{Noah Miller}
\date{March 30, 2026}

\begin{document}
\maketitle

\section{Introduction}
A vector is an arrow pointing in a specific direction with a magnitude. Vague definition, but it will get you through the most of Physics. \\
I will use the boldface notation for vectors(cuz I think it's cooler).
\begin{equation}
    \mathbf{A}=A_x \hat{x} + A_y \hat{y}
\end{equation}
I will represent the magnitude of $\mathbf{A}$ with $|\mathbf{A}|$, where
\begin{equation}
    |\mathbf{A}|^2=A_x^2+A_y^2
\end{equation}
Or, if I feel lazy(which by the way I always do), I will just represent the magnitude by A. 
\\ I will represent a unit vector by $\hat{\mathbf{A}}$, where--
\begin{equation}
    \hat{\mathbf{A}}=\frac{\mathbf{A}}{A}
\end{equation}
I will represent the unit vectors along the cartesian coordinate axes by $\hat{x},\hat{y}$ and $\hat{z}$ along x,y and z axes respectively.

\section{Vector Operations}
To add vectors, just add their components together.
\begin{equation}
    \mathbf{A}+\mathbf{B}=(A_x+B_x)\hat{x}+(A_y+B_y)\hat{y}
\end{equation}
To multiply a vector by a scalar, multiply each of the vector's component with the scalar.
\begin{equation}
    a\mathbf{A}=aA_x\hat{x}+aA_y\hat{y}
\end{equation}
Geometrically, multiplying a vector with a scalar stretches it to a length of the scalar times the previous magnitude of the vector, with no change in the direction.
If the vector is multiplied by -1, it rotates the vector by 180 degrees.\\
To multiply two vectors, well... there are 2 ways to do so. One is the dot product. Another is the cross product.
\\ The dot product of two vectors is defined as
\begin{equation}
    \mathbf{A} \cdot \mathbf{B}=AB\cos\theta
\end{equation}
Where $\theta$ is the angle between the vectors. Geometrically, the dot product of $\mathbf{A}$ and $\mathbf{B}$ is the product of A and the magnitude of $\mathbf{B}$'s projection over $\mathbf{A}$.
\\ Do a bit of math, you will be able to prove that 
\begin{equation}
    \mathbf{A} \cdot \mathbf{B}=A_xB_x+A_yB_y
\end{equation}
The dot product is obviously commutative. \\
The second type of product is the cross product. It is defined as 
\begin{equation}
    \mathbf{A}\times\mathbf{B}=\begin{vmatrix}
                                \hat{x} & \hat{y} & \hat{z} \\
                                A_x     &   A_y   &   A_z   \\
                                B_x     &   B_y   &   B_z
                                \end{vmatrix}
\end{equation} 
\begin{equation}
    |\mathbf{A}\times\mathbf{B}|=AB\sin\theta
\end{equation}
Geometrically, the cross product of $\mathbf{A}$ and $\mathbf{B}$ is a vector of magnitude equal to the area of the parallelogram formed by the two vectors pointing in a direction perpendicular to the plane formed by $\mathbf{A}$ and $\mathbf{B}$.
Of course, there can be two vectors which are perpendicular to the plane formed by the two vectors. Thus, we resort to the right hand thumb rule. I ain't writing the rule. If you forgot it, kys. 
Cross product is not commutative. 

\section{Vector products}
Apart from dot and cross products, there are two other type of products. Namely, Scalar Triple Product and Vector Triple product. \\
The scalar triple product is defined as $\mathbf{A}\cdot(\mathbf{B}\times\mathbf{C})$. Geometrically, this is the volume of the parallelopiped formed by the three vectors.
Note that:-
\begin{equation}
    \mathbf{A}\cdot(\mathbf{B}\times\mathbf{C})=\mathbf{C}\cdot(\mathbf{A}\times\mathbf{B})=\mathbf{B}\cdot(\mathbf{C}\times\mathbf{A})
\end{equation}
i.e, alphabetical order is preserved. If the order is not alphabetical, the magnitude of the triple product will stay the same. But the sign will be opposite.
\\ The vector triple product is defined as $\mathbf{A}\times(\mathbf{B}\times\mathbf{C})$.\\
This can be simplified using the \textbf{BAC-CAB rule}:-
\begin{equation}
    \mathbf{A}\times(\mathbf{B}\times\mathbf{C})=\mathbf{B}(\mathbf{A}\cdot\mathbf{C})-\mathbf{C}(\mathbf{A}\cdot\mathbf{B})
\end{equation}
Proof? Trust me. It works.

\section{Gradient}
To introduce the concept, I will use an example of temperature(which is one of the infinite cases one could think of).
Suppose I have a function T(x,y,z) with tells me the temperature of at the point (x,y,z). Now, I wanna find the change in temperature in moving from
(x,y,z) to (x+dx,y+dy,z+dz). 
\begin{equation}
    dT=T(x+dx,y+dy,z+dz)-T(x,y,z)=\frac{\partial T}{\partial x}dx + \frac{\partial T}{\partial y}dy + \frac{\partial T}{\partial z}dz
\end{equation}
Just because I can, I will write
\begin{align}
    dT = \left(\frac{\partial T}{\partial x}\hat{x} + \frac{\partial T}{\partial y}\hat{y} + \frac{\partial T}{\partial z}\hat{z}\right)\cdot(dx\hspace{0.05cm}\hat{x}+dy\hspace{0.05cm}\hat{y}+dz\hspace{0.05cm}\hat{z})= \nabla T \cdot d\mathbf{r}
\end{align}
Where
\begin{equation}
    \nabla T = \frac{\partial T}{\partial x}\hat{x} + \frac{\partial T}{\partial y}\hat{y} + \frac{\partial T}{\partial z}\hat{z}
\end{equation}
is called the \textbf{Gradient} of T. \\
Geometrically speaking, $\nabla T$ points in the direction where dT is maximum. The magnitude of $\nabla T$ gives the maximum rate of increase per unit distance.
\begin{equation}
    \mathbf{r}=x\hspace{0.05cm}\hat{x}+y\hspace{0.05cm}\hat{y}+z\hspace{0.05cm}\hat{z} \notag 
\end{equation}

\begin{equation}
    d\mathbf{r}=dx\hspace{0.05cm}\hat{x}+dy\hspace{0.05cm}\hat{y}+dz\hspace{0.05cm}\hat{z} \notag
\end{equation}
Should have had defined these two way earlier but nvm. Better late than never/ \\
One could imagine the gradient this way. Suppose you are climbing a hill. You find the steepest ascent. This is the gradient. \\
Let's look at a case where $\nabla T = 0$. This means that the rate of change of temperature over d\textbf{r} there is zero. This could either correspond to a maximum(the top of the hill), a minimum(the bottom of the hill) or a saddle point(a pass).

\section{The Del Operator}
We can write (13) as 
\begin{equation}
    \nabla T = \left(\hat{x}\frac{\partial}{\partial x}+\hat{y}\frac{\partial}{\partial y}+\hat{z}\frac{\partial}{\partial z}\right)T
\end{equation}
As temperature T is a scalar, we could define 
\begin{equation}
    \nabla = \hat{x}\frac{\partial}{\partial x}+\hat{y}\frac{\partial}{\partial y}+\hat{z}\frac{\partial}{\partial z}
\end{equation}
This is the \textbf{Del Operator}. \\ Although it is obviously not a vector, it lowkey behaves like one.
You could take its dot product(divergence) or cross product(curl) with another vector.

\section{Divergence and Curl}
The divergence of a vector \textbf{V} is defined as
\begin{equation}
    \nabla \cdot \mathbf{V} = \frac{\partial V_x}{\partial x}+\frac{\partial V_y}{\partial y}+\frac{\partial V_z}{\partial z}
\end{equation}
It gives us a measure of how much a vector ``spreads out" at a particular point. Note that \textbf{V} is a function. Sketching the graph of \textbf{V}
will give us a \textbf{Vector Field}. \\
The curl of a vector is defined as 
\begin{equation}
    \nabla \times \mathbf{V}=\begin{vmatrix}
                                \hat{x} & \hat{y} & \hat{z} \\
                                \partial/\partial x & \partial/\partial y & \partial/\partial z \\
                                 V_x & V_y & V_z
                              \end{vmatrix}
\end{equation}
The curl of a vector function gives us a measure of how much a vector field ``swirls" around a particular point in space.

\section{Product rules}
There is nothing to explain here. Just a bunch of properties and formulae you gotta memorize. Good luck.
\begin{align}
    \nabla(fg)=f\nabla g + g \nabla f \hspace{4.5cm} \\
    \nabla (\mathbf{A} \cdot \mathbf{B})=\mathbf{A}\times(\nabla \times \mathbf{B})+\mathbf{B}\times(\nabla\times\mathbf{A})+(\mathbf{A}\cdot \nabla)\mathbf{B}+(\mathbf{B}\cdot\nabla)\mathbf{A} \\
    \nabla\cdot(f\mathbf{A})=f(\nabla\cdot\mathbf{A})+\mathbf{A}\cdot(\nabla f)\hspace{2.5cm} \\
    \nabla\cdot(\mathbf{A}\times\mathbf{B})=\mathbf{B}\cdot(\nabla \times \mathbf{A})-\mathbf{A}\cdot(\nabla \times \mathbf{B})\hspace{0.7cm}\\
    \nabla \times (f\mathbf{A})=f(\nabla \times \mathbf{A})-\mathbf{A}\times(\nabla f) \hspace{1.85cm}\\
    \nabla \times (\mathbf{A}\times\mathbf{B})=(\mathbf{B}\cdot \nabla)\mathbf{A} - (\mathbf{A}\cdot \nabla)\mathbf{B} + \mathbf{A}(\nabla\cdot\mathbf{B})-\mathbf{B}(\nabla\cdot\mathbf{A}) \hspace{0.55cm}
\end{align}
f and g are scalar functions. \\
No this ain't Hieroglyphics. This is Vector Calculus. Thank you.
\section{Index notations}
I will define a bunch of fancy notations in this section, startinf off with the \textbf{Levi-Civita symbols}.\\
For that, I will define three sets first. 

\begin{align}
    S_1 = {(1,2,3),(3,1,2),(2,3,1)}\notag\\
    S_2 = {(3,2,1),(2,1,3),(1,3,2)}\notag\\
    S_3 = \{(i,j,k): \text{i=j  or i=k or j=k}\}\notag\\
\end{align}
Now, 
\[
\epsilon_{ijk} = \left\{ 
\begin{array}{lll}
   1 & \text{if } (i,j,k) \in S_1 & \text{+ve} \\
   -1   & \text{if } (i,j,k) \in S_2 & \text{-ve} \\
   0  & \text{if } (i,j,k) \in S_3 & \text{0}
\end{array} 
\right.
\]
In terms of the Levi-Civita symbols, I can write the cross product as 
\begin{equation}
    (\mathbf{u}\times\mathbf{v})_i = \epsilon_{ijk}u_jv_k
\end{equation}
The dot product is just 
\begin{equation}
    \mathbf{u}\cdot{v}=u_iv_i
\end{equation}
You can also write the divergence and curls in the index notations. Just replace \textbf{u} with $\nabla$.
Honestly? So far, I've only seen these notations being used for proof-writing. Probably they have some real uses(although proof-writing is indeed a real use, but.. you get my point.).
\section{Higher derivatives}

There are 5 types of higher derivatives:-\vspace{0.3cm}\\
(i) Divergence of gradient: $\nabla \cdot (\nabla T)$ \\
(ii) Curl of gradient: $\nabla \times (\nabla T)$ \\
(iii) Divergence of curl: $\nabla \cdot (\nabla \times \mathbf{V})$ \\
(iv) Curl of curl: $\nabla \times (\nabla \times \mathbf{V})$ \\
(v) Gradient of divergence: $\nabla(\nabla\cdot \mathbf{V})$ \vspace{0.3cm}\\
If you expand (i), you will get
\begin{equation}
    \nabla \cdot (\nabla T)=\frac{\partial^2T}{\partial x^2}+\frac{\partial^2T}{\partial y^2}+\frac{\partial^2T}{\partial z^2}=\nabla^2 T
\end{equation}
(24) is called the \textbf{Laplacian of T}. \\
(ii) is always zero. \\
(iii) is alwayz zero. \\
(iv) can be expanded as 
\begin{equation}
    \nabla \times (\nabla \times \mathbf{V})=\nabla(\nabla \cdot \mathbf{V})-\nabla^2 \mathbf{V}
\end{equation}

\section{Line Integral}
A line integral is defined as 
\begin{equation}
    I_l = \int_{a}^{b} \mathbf{V}\cdot d\mathbf{l}
\end{equation}
Where \textbf{V} is a vector function, and d\textbf{l} is an infinitesimal displacement vector along the path along which the line integral is evaluated. \\
If the path is closed, i.e, a=b, then $I_l$ will be called a \textbf{closed integral}, which is represented by 
\begin{equation}
    I_l = \oint \mathbf{V}\cdot d\mathbf{l}
\end{equation}
The value of a line integral depends on the path taken, but there is a specific class of vectors for which the line integral doesn't depend on the path.

\section{Surface Integral}
A surface integral over a surface S is defined as
\begin{equation}
    I_s = \int_S \mathbf{V}\cdot d\mathbf{a}
\end{equation}
Where \textbf{V} is again a vector function, and d\textbf{a} is a small patch of area on the surface about which the surface integral is to be calculated. If the surface is closed, we call $I_s$ a \textbf{closed surface integral}, which 
is again represented as 
\begin{equation}
    I_s = \oint \mathbf{V}\cdot d\mathbf{a}
\end{equation}

\section{Volume Integral}
A volume integral is defined as 
\begin{equation}
    I_v=\int_{V}T\hspace{0.01cm}d\tau 
\end{equation}
Where T is a scalar function and $d\tau$ is an infinitesimal volume element. In cartesian coordinates,
\begin{equation}
    d\tau=dx\hspace{0.03cm}dy\hspace{0.03cm}dz
\end{equation}
\section{The Fundamental Theorems of Vector Calculus}
Hopefully you didn't skip high school. So you must know that for a scalar function f,
\begin{equation}
    \int_{a}^{b}\left(\frac{df}{dx}\right)dx=f(b)-f(a)
\end{equation}
In vector calculus, there are three different types of `derivatives'-- gradient, divergence and curl.
Each one has its own fundamental theorem. \\
The \textbf{Fundamental Theorem for Gradients} goes on to say that for a scalar function T(x,y,z) -- 
\begin{equation}
    \int_{a}^{b}\nabla T\cdot d\mathbf{l}=T(b)-T(a)
\end{equation}
Two interesting points to note:--\\
(i)The line integral in (34) does not depend on the path. Turns out that gradients have this interesting property of making their scalar functions path-independent. \\
(ii)$\oint\nabla T\cdot d\mathbf{l}=0$, for obvious reasons.
\\ 
\\
The \textbf{Fundamental Theorem for divergence} goes on to say that--
\begin{equation}
    \int_{V}(\nabla\cdot\mathbf{V})d\tau = \oint_{S}\mathbf{V}\cdot d\mathbf{a}
\end{equation}
For reasons I don't know, (35) has three special names: \textbf{Gauss' Theorem, Green's Theorem} or simply the \textbf{Divergence Theorem}.\\
Geometrically one could think like this: if \textbf{V} represents the flow of an incompressible fluid, the flux of \textbf{V}(the RHS of Eq.35) is the same as the volume flow rate of the fluid through the volume.\\
\\
\textbf{The Fundamental Theorem of Curls} states that --
\begin{equation}
    \int_{S}(\nabla \times \mathbf{V})\cdot d\mathbf{a}=\oint_{P}\mathbf{V}\cdot d\mathbf{l}
\end{equation}
This is also called the \textbf{Stokes' Theorem}. \\
Geometrical intepretation: the curl of a vector field measures the ``swirl'' of the vectors around a point. To get a collective measure of swirl, one could find the flux of the curl along a surface \textit{S}.
Now, imagine a pond, in which the water is rotating in a whirlpool. If you look at the edges of the whirlpool, you can see that water is flowing out of/into the edges of the whirlpool.
The more the whirlpool swirls, the more more water flows out off/into the edges. \\
The \textit{P} here is the perimeter of the surface \textit{S}. \\The RHS of (36) is thus called the ``circulation'' of \textbf{V}. \\
Two interesting points to note:--\\
(i) $\int_{S}(\nabla \times \mathbf{V})\cdot d\mathbf{a}$ depends only on the boundary line, and not the area of the surface as such. \\
(ii)$\oint_{S}(\nabla \times \mathbf{V})\cdot d\mathbf{a} = 0$, as the boundary line will shrink to a point for a closed surface, just like a balloon. \\

\section{Integration by Parts}
Let's first look at how we integrated by parts for scalar functions. 
\begin{equation}
    \frac{d}{dx}(fg)=f\left(\frac{dg}{dx}\right)+g\left(\frac{df}{dx}\right)
\end{equation}
Integrating LHS and RHS from a to b --
\begin{align}
\int_{a}^{b}\frac{d}{dx}(fg)dx = \int_{a}^{b}f\left(\frac{dg}{dx}\right)dx + \int_{a}^{b}g\left(\frac{df}{dx}\right)dx\notag \\
\implies \int_{a}^{b}f\left(\frac{dg}{dx}\right)dx = \left. (fg)\right|_{a}^{b}-\int_{a}^{b}g\left(\frac{df}{dx}\right)dx
\end{align}
We proceed in a similar fashion while integrating vector function by parts. The product rule for gradients is-- 
\begin{equation}
    \nabla\cdot(f\mathbf{A}) = f(\nabla\cdot\mathbf{A})+A\cdot(\nabla f) \notag
\end{equation}
Integrating the LHS and RHS over an infinitesimal volume $d\tau$ and invoking Green's theorem(I like this name),
\begin{align}
    \int \nabla\cdot(f\mathbf{A}) d\tau = \int f(\nabla\cdot\mathbf{A})d\tau+\int A\cdot(\nabla f)d\tau = \oint \nabla f\mathbf{A}\cdot d\mathbf{a} \notag \\
    \implies \int_{V} f(\nabla\cdot\mathbf{A})d\tau = \oint_{S} f\mathbf{A}\cdot d\mathbf{a}-\int_{V} A\cdot(\nabla f)d\tau \hspace{2.8cm}
\end{align}

\section{Spherical Coordinates}
As the name suggests, we write the coordintes with reference to a sphere. As in, we imagine the particle moving along the radii of one or more concentric
spheres. \\
In cartesian coordinates, we have three perpendicular coordinate axes for reference. To specify the position of the particle, we write the displacement of the particle along these axes.
For example, if I say that the coordinates of a particle are (2,3,7) in the cartesian system, it means that the particle is 2 units away from the origin along the x-axis, 3 units away from the 
origin along the y-axis and 7 units away from the origin along the z-axis. \\
In Spherical Coordinates, we specify the position of the particle by writing its distance from the origin(r), the angle the position vector of the particle makes with the z-axis, or simply the \textbf{Plane angle}($\theta$), and the
angle made by the position vector's projection with the x-axis, or simply the \textbf{Azimuthal Angle}($\phi$). \\
The unique position of the particle in the three dimensional space is denoted by (r,$\theta$,$\phi$). \\
We define three unit vectors:-\\
(i)$\hat{r}$ along \textbf{r}. \\
(ii)$\hat{\theta}$ perpendicular to $\hat{r}$, in the direction where $\theta$ is increasing. \\ 
(iii)$\hat{\phi}$ perpendicular to both $\hat{r}$ and $\hat{\theta}$, in the direction where $\phi$ is increasing. \\ \\
Thus, one can write a vector using spherical coordinates as--
\begin{equation}
    \mathbf{A}=A_r\hat{r}+A_{\theta}\hat{\theta}+A_{\phi}\hat{\phi}
\end{equation}
The transformation equations are--
\begin{align}
    z = r\cos\theta \hspace{0.9cm}\\
    x = r\sin\theta\cos\phi \\
    y = r\sin\theta\sin\phi \\ \notag
\end{align}
For unit vectors--
\begin{align}
    \hat{r} =\sin\theta\cos\phi\hspace{0.1cm}\hat{x}+\sin\theta\sin\phi\hspace{0.1cm}\hat{y}+\cos\theta\hspace{0.1cm}\hat{z} \\
    \hat{\theta} =\cos\theta\cos\phi\hspace{0.1cm}\hat{x}+\cos\theta\sin\phi\hspace{0.1cm}\hat{y}-\sin\theta\hspace{0.1cm}\hat{z} \\
    \hat{\phi} = -\sin\phi\hspace{0.1cm}\hat{x}+\cos\phi\hspace{0.1cm}\hat{y}\hspace{3cm} \\ \notag
\end{align}
While writing the displacement vector $d\mathbf{l}$ in spherical coordinates--
\begin{align}
    d\mathbf{l}= dl_{r}\hat{r}+dl_{\theta}\hat{\theta}+dl_{\phi}\hat{\phi} = dr\hspace{0.03cm}\hat{r}+ r\hspace{0.03cm}d\theta\hspace{0.1cm}\hat{\theta}+r\sin\theta d\phi \hspace{0.1cm}\hat{\phi}
\end{align} 
To write $d\tau$ in spherical coordinates:- 
\begin{align}
    d\tau=r^2 \sin\theta dr\hspace{0.05cm}d\theta\hspace{0.05cm}d\phi
\end{align}
I will write the expressions for gradient, divergence, curl and laplacian in spherical coordinates. See Appendix A of Griffiths' \textit{Introduction to Electrodynamics} for the proof.
\begin{align}
    \nabla T = \frac{\partial T}{\partial r}\hat{r}+\frac{1}{r}\frac{\partial T}{\partial \theta}\hat{\theta}+\frac{1}{r\sin\theta}\frac{\partial T}{\partial \phi}\hat{\phi} \hspace{10cm} \\
    \nabla \cdot \mathbf{v} = \frac{1}{r^2}\frac{\partial}{\partial r}(r^2v_r)+\frac{1}{r\sin\theta}\frac{\partial}{\partial \theta}(v_{\theta}\sin\theta)+\frac{1}{r\sin\theta}\frac{\partial v_{\phi}}{\partial\phi}\hspace{6.6cm} \\
    \nabla \times \mathbf{v} = \frac{1}{r\sin\theta}\left[\frac{\partial}{\partial \theta}(v_{\phi}\sin\theta)-\frac{\partial v_{\theta}}{\partial \phi}\right]\hat{r}+\frac{1}{r}\left[\frac{1}{\sin\theta}\frac{\partial v_r}{\partial \phi}-\frac{\partial}{\partial r}(rv_{\phi})\right]\hat{\theta} + \frac{1}{r}\left[\frac{\partial}{\partial r}(rv_{\theta})-\frac{\partial v_r}{\partial\theta}\right]\hat{\phi} \\
    \nabla^2T=\frac{1}{r^2}\frac{\partial}{\partial r}\left(r^2 \frac{\partial T}{\partial r}\right)+\frac{1}{r^2 \sin\theta}\frac{\partial}{\partial\theta}\left(\sin\theta \frac{\partial T}{\partial \theta}\right)+\frac{1}{r^2\sin^2\theta}\frac{\partial^2T}{\partial\phi^2} \hspace{5cm}
\end{align}
Fucking Christ I had to fry 3 gazillion braincells just to type this shit out...

\section{Cylindrical Coordinates}
Here, instead of a sphere or rectancle, we take a cylinder as a reference. We specify the position of a particle using 3 parameters --\\
(i) Distance from the axis of the cylinder(s)\\
(ii) Azimuthal angle($\phi$)\\
(iii) Displacement along the height of the cylinder(z) \\
\\
Their relation to cartesian coordinates is --
\begin{align}
    x = s\cos\phi \\
    y = s\sin\phi \\
    z = z \hspace{0.9cm}
\end{align}
Thus, 
\begin{equation}
    d\mathbf{l}=ds\hspace{0.1cm}\hat{s}+s\hspace{0.1cm}d\phi\hspace{0.1cm}\hat{\phi}+dz\hspace{0.1cm}\hat{z}
\end{equation}
The volume element will be --
\begin{equation}
    d\tau = s\hspace{0.1cm}ds\hspace{0.05cm}d\phi\hspace{0.05cm}dz
\end{equation}
The expressions of gradient, divergence and curl in cylindrical coordinates are(LaTeX torture again..) ---
\begin{align}
    \nabla T = \frac{\partial T}{\partial s}\hat{s}+\frac{1}{s}\frac{\partial T}{\partial \phi}\hat{\phi}+\frac{\partial T}{\partial z}\hat{z}\hspace{7cm}\\
    \nabla\cdot\mathbf{v} = \frac{1}{s}\frac{\partial}{\partial s}(sv_{s})+\frac{1}{s}\frac{\partial v_{\phi}}{\partial \phi}+\frac{\partial v_z}{\partial z}\hspace{6.2cm}\\
    \nabla\times\mathbf{v} = \left(\frac{1}{s}\frac{\partial v_z}{\partial \phi}-\frac{\partial v_{\phi}}{\partial z}\right)\hat{s}+\left(\frac{\partial v_s}{\partial z}-\frac{\partial v_z}{\partial s}\right)\hat{\phi}+\frac{1}{s}\left[\frac{\partial}{\partial s}(sv_{\phi})-\frac{\partial v_s}{\partial \phi}\right]\hat{z}\\
    \nabla^{2}T = \frac{1}{s}\frac{\partial}{\partial s}\left(s\frac{\partial T}{\partial s}\right)+\frac{1}{s^2}\frac{\partial^2 T}{\partial \phi^2}+\frac{\partial^2 T}{\partial z^2}\hspace{5.3cm}
\end{align}

\section{One Dimensional Dirac Delta Function}
Before studying the Dirac Delta function, let's look at the function---
\begin{equation}
    \mathbf{v}=\frac{1}{r^2}\hat{r}
\end{equation}
We can compute its divergence using (49).
\begin{equation}
    \nabla\cdot\mathbf{v}=\frac{1}{r^2}\frac{\partial}{\partial r}(1)=0
\end{equation}
Now, let's find the flux of $\mathbf{v}$ over a sphere of radius R.

\begin{equation}
    \oint \mathbf{v}\cdot d\mathbf{a}=\int_{0}^{2\pi}\int_{0}^{\pi} \left(\frac{1}{R^2}\hat{r}\right)\cdot R^2 \sin\theta d \theta d\phi\hspace{0.1cm}\hat{r} = 4\pi
\end{equation}
But, if we compute (64) using the Divergence theorem,
\begin{equation}
    \oint \mathbf{v}\cdot d\mathbf{a} = \int_{V}(\nabla\cdot\mathbf{v})d\tau=0
\end{equation}
Of course (64) and (65) can't be true simultaneously. So, which is correct? And more importantly, is the divergence theorem wrong?\\
Turns out that all discrepancies lie within the point r=0. The function blows up at r=0. In computing the integral, we are essentially dividing by zero at the point r=0. \\
The divergence is indeed zero, \textit{except} for the point where r=0. The Dirac Delta function will help us overcome this discrepancy.\\\\
The Dirac Delta function is defined as:--
\begin{equation}
    \delta(x)=\begin{cases}
        0 & \text{if }x\ne 0 \\
        \infty & \text{if }x=0
    \end{cases}
\end{equation}
and
\begin{equation}
    \int_{-\infty}^{\infty}\delta(x)dx=1
\end{equation}
Now suppose we have a smooth and continuous function(for safety reasons) $f(x)$. I take its product with the dirac delta function.\\
$f(x)\delta(x)$ is zero everywhere \textit{except} at x=0.\\
Thus, we can write --
\begin{equation}
    f(x)\delta(x)=f(0)\delta(x)
\end{equation}
Now, I will integrate this product from $-\infty$ to $+\infty$:--
\begin{equation}
    \int_{-\infty}^{\infty}f(x)\delta(x)dx=f(0)\int_{-\infty}^{\infty}\delta(x)dx=f(0)
\end{equation}
If you've studied analytical geometry in High school(which I suppose you have), you know that we can shift the origin fron (0,0) to (a,0). In that case, all the x's will be replaced by (x-a)'s. 
\begin{equation}
    f(x-a)\delta(x-a)=f(a)\delta(x-a)
\end{equation}
\begin{equation}
    \int_{-\infty}^{\infty}f(x-a)\delta(x-a)dx=f(a)\int_{-\infty}^{\infty}\delta(x-a)dx=f(a)
\end{equation}
Note that the integral was not affected `cause I just shifted the plot rightwards by `a' units. Thus, the area under the curve stays the same. \\
Another interesting property of the dirac delta function is:--
\begin{equation}
    \delta(kx)=\frac{1}{|k|}\delta(x)
\end{equation}
To prove this, I will drag $\delta(kx)$ into an integral:--
\begin{equation}
    \int_{-\infty}^{\infty}f(x)\delta(kx)dx
\end{equation}
Now, I will make a change in variable.
\begin{align}
    y=kx
\end{align}
Here's the catch. If k is positive, chef's kiss. But if not, the $-\infty$ and $\infty$ in (73) will reverse orders. To account for this flipping of the infinites, a minus has to be included.
Thus, (73) can be re-written as:--
\begin{equation}
    \pm\frac{1}{k}\int_{-\infty}^{\infty}f\left(\frac{y}{k}\right)\delta(y)dy=\frac{f(0)}{|k|}
\end{equation}
Comparing this with (69), we can say that (72) is true.

\section{Three Dimensional Dirac Delta Function}
The delta function can be generalized to three dimensions as--
\begin{equation}
    \delta^3(\mathbf{r}) =\delta(x)\delta(y)\delta(z)
\end{equation}
$\delta^3(\mathbf{r})$ is zero everywhere, save for the point (0,0,0).\\
The volume integral of the delta function over the 3 dimensional space is 1.
\begin{equation}
    \int_{V}\delta(x)\delta(y)\delta(z)d\tau=\int_{-\infty}^{\infty}\int_{-\infty}^{\infty}\int_{-\infty}^{\infty}\delta(x)\delta(y)\delta(z)dx\hspace{0.1cm}dy\hspace{0.1cm}dz=1
\end{equation} 
Also,
\begin{equation}
    \int_{-\infty}^{\infty}f(\mathbf{a})\delta^3(\mathbf{r}-\mathbf{a})d\tau=f(\mathbf{a})
\end{equation}
We are now in a position to resolve the discrepancy which I told you of at the beginning of section 16. 
As we saw, $\nabla\cdot\mathbf{v}$ was zero everywhere, but blew up at (0,0,0). And its integral over $d\tau$ was $4\pi$. Keeping all these conditions in mind, we can say--
\begin{equation}
    \nabla\cdot\mathbf{\frac{1}{r^2}\hat{r}}=4\pi\delta^3(\mathbf{r})
\end{equation}

\section{Helmholtz Theorem and Vector Fields}
The Helmholtz theorem states that if the divergence and curl of a vector function \textbf{F} go to zero faster than $1/r^2$ as $r\rightarrow \infty$, and \textbf{F}(\textbf{r}) goes to zero as $\mathbf{r}\rightarrow \infty$, then \textbf{F} is uniquely given by:--
\begin{equation}
    \mathbf{F}=-\nabla U + \nabla\times\mathbf{W}
\end{equation}
See \textit{Appendix B} of Griffiths' \textit{Introduction to Electrodynamics} for the proof.
\end{document}
